% Options for packages loaded elsewhere
\PassOptionsToPackage{unicode}{hyperref}
\PassOptionsToPackage{hyphens}{url}
\PassOptionsToPackage{dvipsnames,svgnames,x11names}{xcolor}
%
\documentclass[
  letterpaper,
  DIV=11,
  numbers=noendperiod,
  oneside]{scrartcl}

\usepackage{amsmath,amssymb}
\usepackage{lmodern}
\usepackage{iftex}
\ifPDFTeX
  \usepackage[T1]{fontenc}
  \usepackage[utf8]{inputenc}
  \usepackage{textcomp} % provide euro and other symbols
\else % if luatex or xetex
  \usepackage{unicode-math}
  \defaultfontfeatures{Scale=MatchLowercase}
  \defaultfontfeatures[\rmfamily]{Ligatures=TeX,Scale=1}
\fi
% Use upquote if available, for straight quotes in verbatim environments
\IfFileExists{upquote.sty}{\usepackage{upquote}}{}
\IfFileExists{microtype.sty}{% use microtype if available
  \usepackage[]{microtype}
  \UseMicrotypeSet[protrusion]{basicmath} % disable protrusion for tt fonts
}{}
\makeatletter
\@ifundefined{KOMAClassName}{% if non-KOMA class
  \IfFileExists{parskip.sty}{%
    \usepackage{parskip}
  }{% else
    \setlength{\parindent}{0pt}
    \setlength{\parskip}{6pt plus 2pt minus 1pt}}
}{% if KOMA class
  \KOMAoptions{parskip=half}}
\makeatother
\usepackage{xcolor}
\usepackage[top=30mm,bottom=30mm,left=15mm,right=80mm,marginparwidth=60mm,marginparsep=10mm,heightrounded]{geometry}
\setlength{\emergencystretch}{3em} % prevent overfull lines
\setcounter{secnumdepth}{-\maxdimen} % remove section numbering
% Make \paragraph and \subparagraph free-standing
\ifx\paragraph\undefined\else
  \let\oldparagraph\paragraph
  \renewcommand{\paragraph}[1]{\oldparagraph{#1}\mbox{}}
\fi
\ifx\subparagraph\undefined\else
  \let\oldsubparagraph\subparagraph
  \renewcommand{\subparagraph}[1]{\oldsubparagraph{#1}\mbox{}}
\fi


\providecommand{\tightlist}{%
  \setlength{\itemsep}{0pt}\setlength{\parskip}{0pt}}\usepackage{longtable,booktabs,array}
\usepackage{calc} % for calculating minipage widths
% Correct order of tables after \paragraph or \subparagraph
\usepackage{etoolbox}
\makeatletter
\patchcmd\longtable{\par}{\if@noskipsec\mbox{}\fi\par}{}{}
\makeatother
% Allow footnotes in longtable head/foot
\IfFileExists{footnotehyper.sty}{\usepackage{footnotehyper}}{\usepackage{footnote}}
\makesavenoteenv{longtable}
\usepackage{graphicx}
\makeatletter
\def\maxwidth{\ifdim\Gin@nat@width>\linewidth\linewidth\else\Gin@nat@width\fi}
\def\maxheight{\ifdim\Gin@nat@height>\textheight\textheight\else\Gin@nat@height\fi}
\makeatother
% Scale images if necessary, so that they will not overflow the page
% margins by default, and it is still possible to overwrite the defaults
% using explicit options in \includegraphics[width, height, ...]{}
\setkeys{Gin}{width=\maxwidth,height=\maxheight,keepaspectratio}
% Set default figure placement to htbp
\makeatletter
\def\fps@figure{htbp}
\makeatother

\usepackage{lastpage}
\usepackage{fancyhdr}
\usepackage{titling}
\pagestyle{fancy}
\fancyhead[L]{CMPE2150 Project 01}
\fancyhead[R]{Semiconductor Control}
\fancyfoot[C]{\thepage/\pageref*{LastPage}}
\setlength{\droptitle}{-2cm}
\KOMAoption{captions}{tablesignature}
\makeatletter
\@ifpackageloaded{tcolorbox}{}{\usepackage[many]{tcolorbox}}
\@ifpackageloaded{fontawesome5}{}{\usepackage{fontawesome5}}
\definecolor{quarto-callout-color}{HTML}{909090}
\definecolor{quarto-callout-note-color}{HTML}{0758E5}
\definecolor{quarto-callout-important-color}{HTML}{CC1914}
\definecolor{quarto-callout-warning-color}{HTML}{EB9113}
\definecolor{quarto-callout-tip-color}{HTML}{00A047}
\definecolor{quarto-callout-caution-color}{HTML}{FC5300}
\definecolor{quarto-callout-color-frame}{HTML}{acacac}
\definecolor{quarto-callout-note-color-frame}{HTML}{4582ec}
\definecolor{quarto-callout-important-color-frame}{HTML}{d9534f}
\definecolor{quarto-callout-warning-color-frame}{HTML}{f0ad4e}
\definecolor{quarto-callout-tip-color-frame}{HTML}{02b875}
\definecolor{quarto-callout-caution-color-frame}{HTML}{fd7e14}
\makeatother
\makeatletter
\makeatother
\makeatletter
\makeatother
\makeatletter
\@ifpackageloaded{caption}{}{\usepackage{caption}}
\AtBeginDocument{%
\ifdefined\contentsname
  \renewcommand*\contentsname{Table of contents}
\else
  \newcommand\contentsname{Table of contents}
\fi
\ifdefined\listfigurename
  \renewcommand*\listfigurename{List of Figures}
\else
  \newcommand\listfigurename{List of Figures}
\fi
\ifdefined\listtablename
  \renewcommand*\listtablename{List of Tables}
\else
  \newcommand\listtablename{List of Tables}
\fi
\ifdefined\figurename
  \renewcommand*\figurename{Figure}
\else
  \newcommand\figurename{Figure}
\fi
\ifdefined\tablename
  \renewcommand*\tablename{Table}
\else
  \newcommand\tablename{Table}
\fi
}
\@ifpackageloaded{float}{}{\usepackage{float}}
\floatstyle{ruled}
\@ifundefined{c@chapter}{\newfloat{codelisting}{h}{lop}}{\newfloat{codelisting}{h}{lop}[chapter]}
\floatname{codelisting}{Listing}
\newcommand*\listoflistings{\listof{codelisting}{List of Listings}}
\makeatother
\makeatletter
\@ifpackageloaded{caption}{}{\usepackage{caption}}
\@ifpackageloaded{subcaption}{}{\usepackage{subcaption}}
\makeatother
\makeatletter
\@ifpackageloaded{tcolorbox}{}{\usepackage[many]{tcolorbox}}
\makeatother
\makeatletter
\@ifundefined{shadecolor}{\definecolor{shadecolor}{rgb}{.97, .97, .97}}
\makeatother
\makeatletter
\@ifpackageloaded{sidenotes}{}{\usepackage{sidenotes}}
\@ifpackageloaded{marginnote}{}{\usepackage{marginnote}}
\makeatother
\makeatletter
\makeatother
\ifLuaTeX
  \usepackage{selnolig}  % disable illegal ligatures
\fi
\IfFileExists{bookmark.sty}{\usepackage{bookmark}}{\usepackage{hyperref}}
\IfFileExists{xurl.sty}{\usepackage{xurl}}{} % add URL line breaks if available
\urlstyle{same} % disable monospaced font for URLs
\hypersetup{
  pdftitle={CMPE2150 Project 01},
  pdfauthor={AJ Armstrong},
  colorlinks=true,
  linkcolor={blue},
  filecolor={Maroon},
  citecolor={Blue},
  urlcolor={Blue},
  pdfcreator={LaTeX via pandoc}}

\title{CMPE2150 Project 01}
\usepackage{etoolbox}
\makeatletter
\providecommand{\subtitle}[1]{% add subtitle to \maketitle
  \apptocmd{\@title}{\par {\large #1 \par}}{}{}
}
\makeatother
\subtitle{Semiconductor Control}
\author{AJ Armstrong}
\date{01 Oct 2024}

\begin{document}
\maketitle
\ifdefined\Shaded\renewenvironment{Shaded}{\begin{tcolorbox}[boxrule=0pt, interior hidden, borderline west={3pt}{0pt}{shadecolor}, frame hidden, enhanced, breakable, sharp corners]}{\end{tcolorbox}}\fi

\renewcommand*\contentsname{Table of contents}
{
\hypersetup{linkcolor=}
\setcounter{tocdepth}{3}
\tableofcontents
}
In this Project, you will design two simple circuits using a
semiconductor to control an isolated DC load (simulated by an LED): by
controlling a electromechanical relay with a BJT and using a
FET-controlled relay.

\begin{tcolorbox}[enhanced jigsaw, toprule=.15mm, titlerule=0mm, colframe=quarto-callout-important-color-frame, colbacktitle=quarto-callout-important-color!10!white, title=\textcolor{quarto-callout-important-color}{\faExclamation}\hspace{0.5em}{Important}, opacitybacktitle=0.6, bottomtitle=1mm, colback=white, leftrule=.75mm, coltitle=black, toptitle=1mm, arc=.35mm, opacityback=0, breakable, rightrule=.15mm, bottomrule=.15mm, left=2mm]

It is important that you show your work in the design sections below.
Make it clear why you selected the values that you did for your design.
Make sure your schematic is neat and clear.

\end{tcolorbox}

\newpage{}

\hypertarget{part-1-bjt-controlled-led}{%
\subsubsection{Part 1: BJT-Controlled
LED}\label{part-1-bjt-controlled-led}}

Design a circuit to control an LED using a BJT:

\begin{itemize}
\item
  On the input side, use a TTL signal and a BJT to drive an LED. The LED
  should be turned On or Off by a \(5V\) or \(0V\) TTL input. You may
  use either a switch or a low-frequency TTL square-wave signal to
  toggle the LED state.
\item
  Make sure the transistor is biased properly, so that in saturation the
  collector current is enough to drive the LED (see below), and that
  \(R_B\) is correct to drive the transistor into saturation for the
  worst case \(\beta\) or \(h_{fe}\) for your transistor. Remember the
  rules of thumb for selecting collector and base resistors in a basic
  digital (cutoff/saturation) BJT biasing.
\item
  On the Collector/Emitter side, configure it so that when the BJT is
  on, an LED is illuminated. Use a \(5V\) supply to power the LED and
  ensure you have calculated an appropriate current-limiting resistor
  \(R_C\). Recall that you can find out the current you need from your
  LED's datasheet. About \(75\%\) of the maximum current is a good
  benchmark if there isn't an optimal current in the spec.
\end{itemize}

For checkoff, attach your schematic and calculations to this document
and demonstrate your circuit for your instructor to sign off
\_\_\_\_\_\_\_\_\_\_\_\_\_\_\_\_\_\_\_\_

\newpage{}

\hypertarget{part-2-fet-controlled-relay}{%
\subsubsection{Part 2: FET-Controlled
Relay}\label{part-2-fet-controlled-relay}}

Design a new circuit that replaces the BJT with a FET-controlled relay:

\begin{itemize}
\item
  On the input side, use a FET to trip your G2R-2 relay. The relay
  should be powered by a \(12 V\) source, and the FET should be turned
  On or Off by a \(5V\) or \(0V\) TTL input. You may use either a switch
  or a low-frequency TTL square-wave signal to toggle the FET state.
\item
  Make sure you have the proper configuration to protect the transistor
  from the electromagnet flyback and its gate from inrush currents.
\item
  The output side will be similar to the previous circuit. Configure it
  so that when the relay is tripped (the FET is on), an LED is
  illuminated. Use a \(5V\) supply to power the LED and ensure you have
  calculated an appropriate current-limiting resistor \(R_L\). Recall
  that you can find out the current you need from your LED's datasheet.
  About \(75\%\) of the maximum current is a good benchmark if there
  isn't an optimal current in the spec.
\end{itemize}

For checkoff, attach your schematic and calculations to this document
and demonstrate your circuit for your instructor to sign off
\_\_\_\_\_\_\_\_\_\_\_\_\_\_\_\_\_\_\_\_



\end{document}
